\subsection{Related Work}

The landscape of efficient pairing verification in constraint systems has evolved through several key contributions. El Housni's \textbf{Pairings in Rank-1 Constraint Systems}~\cite{PR1CS22} established the theoretical framework for optimizing pairings in Rank-1 Constraint Systems, introducing techniques for efficient Miller loop implementation and final exponentiation that exploit the unique properties of constraint-based computation models. This work demonstrated that operations traditionally considered expensive in native implementations could be restructured to achieve remarkable efficiency gains in arithmetized environments.

Feltroidprime's subsequent research on \textbf{enhanced extension field multiplication techniques}~\cite{FXFM23} introduced polynomial verification methods using the Schwartz-Zippel lemma for $\Fe{12}$ arithmetic operations. By representing extension field elements as polynomials and verifying their operations through polynomial arithmetic, this approach achieved significant performance improvements in the pairing verification.

Novakovic and Eagen's work \textbf{On Proving Pairings}~\cite{OPP24} presented breakthrough optimization strategies that enable the elimination of final exponentiation overhead through residue witness techniques. Their approach demonstrates that two elements in $\Fe{12}$ can be proven equivalent without expensive final exponentiation by introducing witness elements that satisfy specific algebraic relationships, and effectively embedding the final exponentiation verification into the Miller loop computation itself.
